%!TEX program = lualatex
% Relatório: Experimento 5 — Transformador
\documentclass[a4paper,12pt]{article}

% =========================================================
% Idioma e tipografia
% =========================================================

\usepackage[brazilian]{babel}
\usepackage{csquotes}
\usepackage{fontspec}
\setmainfont{Latin Modern Roman}


% =========================================================
% Layout
% =========================================================
\usepackage[top=3cm,bottom=2.5cm,left=3cm,right=2.5cm]{geometry}
\usepackage{setspace}
\onehalfspacing
\usepackage{indentfirst}
\setlength{\parindent}{1.25cm}

% =========================================================
% Matemática, tabelas, listas e figuras
% =========================================================
\usepackage{amsmath,amssymb}
\usepackage{float}
\usepackage{booktabs}
\usepackage{array}
\usepackage{enumitem}
\usepackage{graphicx}
\usepackage{xcolor}
\usepackage{listings}

% =========================================================
% Referências bibliográficas
% =========================================================
\usepackage[
  backend=biber,
  style=authoryear,
  sorting=nyt,
  maxcitenames=2
]{biblatex}
\addbibresource{referencias_transformador.bib}

\setlist[itemize]{noitemsep,topsep=3pt}
\setlist[enumerate]{noitemsep,topsep=3pt}
\graphicspath{{./}{./Figs/}}

\definecolor{codigocinza}{RGB}{247,247,247}
\definecolor{codigoverde}{RGB}{0,100,0}
\lstset{
  language=Python,
  basicstyle=\ttfamily\footnotesize,
  backgroundcolor=\color{codigocinza},
  commentstyle=\color{codigoverde},
  keywordstyle=\color{blue!70!black},
  stringstyle=\color{red!65!black},
  numbers=left,
  numberstyle=\tiny\color{gray},
  numbersep=8pt,
  frame=single,
  rulecolor=\color{gray!55},
  breaklines=true,
  breakatwhitespace=false,
  columns=fullflexible,
  keepspaces=true,
  showstringspaces=false,
  tabsize=4
}
\renewcommand{\lstlistingname}{Código}

\newcommand{\imagemrelatorio}[2][width=.7\linewidth]{%
  \IfFileExists{#2}{%
    \includegraphics[#1]{#2}%
  }{%
    \fbox{\parbox[c][3.5cm][c]{.78\linewidth}{%
      \centering Imagem não encontrada: \texttt{\detokenize{#2}}%
    }}%
  }%
}

% =========================================================
% Links
% =========================================================
\usepackage{xurl}
\usepackage{hyperref}
\hypersetup{
  colorlinks=true,
  linkcolor=black,
  urlcolor=blue,
  citecolor=black,
  pdfauthor={Miguel Gayani, José Mauro Elsas e Israel Oliveira Carneiro},
  pdftitle={Experimento 5 — Transformador}
}

% =========================================================
% Dados da capa
% =========================================================
\newcommand{\Instituicao}{Universidade do Estado do Rio de Janeiro}
\newcommand{\Curso}{Física VII -- Oceanografia}
\newcommand{\Professor}{Professor: Alexandre Magnus Gomes Carvalho}
\newcommand{\Alunos}{Alunos: Miguel Gayani, \\
                             Jose Mauro Elsas e \\
                             Israel Oliveira Carneiro}
\newcommand{\Cidade}{Rio de Janeiro}
\newcommand{\Ano}{2026}

\begin{document}
\pagenumbering{gobble}

% =========================================================
% Capa
% =========================================================
\begin{titlepage}
  \centering
  {\large \Instituicao\par}
  \vspace{0.5cm}
  {\large \Curso\par}
  \vspace{1cm}
  {\LARGE\textbf{Experimento 5}\par}
  \vspace{0.6cm}
  {\Large Lei de Faraday e as relações\\
  entre número de espiras e tensão
  \par}
  \vspace{1cm}
  \imagemrelatorio[width=.3\linewidth]{Figs/uerj.jpg}
  \vfill
  \begin{flushright}
    \Alunos\\
    \Professor
  \end{flushright}
  \vspace{1cm}
  {\Cidade{} --- \Ano\par}
\end{titlepage}

\clearpage
\pagenumbering{arabic}

% =========================================================
% Objetivo
% =========================================================
\section{Objetivo}

Verificar experimentalmente a relação entre as tensões eficazes nos enrolamentos
primário e secundário de um transformador e seus respectivos números de espiras.
Para isso, determinar as razões $V_s/V_p$ em duas configurações de alimentação,
calcular suas médias e seus desvios padrão e comparar os resultados com as razões
previstas pelo modelo de transformador ideal.

% =========================================================
% Introdução e fundamentação teórica
% =========================================================
\section{Introdução}

O desenvolvimento dos transformadores está ligado aos estudos sobre a relação
entre eletricidade e magnetismo realizados no século XIX. Em 1831, Michael
Faraday demonstrou a indução eletromagnética: uma variação do fluxo magnético
através de um circuito produz nele uma força eletromotriz induzida. Esse princípio
se tornou a base de geradores, motores e transformadores
\parencite{faraday1832,halliday}.

Com a posterior expansão dos sistemas de corrente alternada, o transformador
passou a desempenhar uma função essencial na transmissão e na distribuição de
energia elétrica. Ele permite elevar a tensão para reduzir a corrente e as perdas
por efeito Joule durante o transporte de energia e, posteriormente, reduzir a
tensão para níveis adequados ao consumo. O mesmo princípio é empregado em fontes
de alimentação, instrumentos eletrônicos, sistemas de isolamento elétrico e
diversos equipamentos de medição e controle.

Um transformador simples é constituído por dois enrolamentos acoplados por um
núcleo ferromagnético. O enrolamento conectado à fonte é denominado
\emph{primário}, independentemente de seu número de espiras; o outro é denominado
\emph{secundário}. Uma tensão alternada aplicada ao primário produz uma corrente
variável e, consequentemente, um fluxo magnético variável no núcleo. Esse fluxo
atravessa o secundário e induz nele uma tensão alternada
\parencite{ling2016,halliday}.

De acordo com a lei de Faraday, a força eletromotriz induzida em um enrolamento
com $N$ espiras é

\begin{equation}
  \mathcal{E}=-N\frac{d\Phi_B}{dt},
  \label{eq:faraday}
\end{equation}

em que $\Phi_B$ é o fluxo magnético através de uma espira. No transformador
ideal, o mesmo fluxo atravessa os dois enrolamentos. Assim, as tensões eficazes
são proporcionais aos respectivos números de espiras:

\begin{equation}
  \frac{V_s}{V_p}=\frac{N_s}{N_p},
  \label{eq:razao_tensoes}
\end{equation}

na qual $V_p$ e $V_s$ são as tensões no primário e no secundário, enquanto
$N_p$ e $N_s$ representam seus números de espiras. Se $N_s>N_p$, o equipamento
é elevador de tensão; se $N_s<N_p$, é abaixador.

Para um transformador ideal, a potência de entrada é igual à potência de saída:

\begin{equation}
  V_p I_p=V_s I_s.
  \label{eq:potencia}
\end{equation}

Combinando as Equações~\eqref{eq:razao_tensoes} e~\eqref{eq:potencia}, obtém-se
a relação inversa entre as correntes:

\begin{equation}
  \frac{I_s}{I_p}=\frac{N_p}{N_s}.
  \label{eq:razao_correntes}
\end{equation}

Transformadores reais apresentam resistência elétrica nos enrolamentos, fluxo de
dispersão, perdas por correntes parasitas e histerese no núcleo, além de
acoplamento magnético imperfeito. Por isso, a razão experimental das tensões pode
diferir da razão entre os números de espiras, sobretudo em equipamentos didáticos
com núcleo móvel ou com entreferros.

% =========================================================
% Materiais e métodos
% =========================================================
\section{Materiais e Métodos}

\subsection{Materiais}

Foram utilizados os seguintes materiais:

\begin{enumerate}[label=\alph*.]
  \item transformador com núcleo ferromagnético laminado e enrolamentos
  de 300 e 1200 espiras;
  \item fonte de tensão alternada com posições de 1 a 10 de ajuste;
  \item dois multímetros digitais iguais, modelo Minipa ET-1110A, utilizados na
  função de voltímetro para corrente alternada;
  \item cabos e pontas de prova para as conexões elétricas.
\end{enumerate}

\begin{figure}[H]
  \centering
  \imagemrelatorio[width=.6\linewidth]{Figs/WhatsApp Image 2026-09-03 at 22.13.04.jpeg}
  \caption{Fonte de Alimentação}
  \label{fig:Fonte}
\end{figure}

\begin{figure}[H]
  \centering
  \imagemrelatorio[width=.6\linewidth]{Figs/transformador.jpeg}
  \caption{transformador utilizado no experimento, com dois
  enrolamentos acoplados por um núcleo ferromagnético laminado.}
  \label{fig:transformador}
\end{figure}

\begin{figure}[H]
  \centering
  \imagemrelatorio[width=.52\linewidth]{Figs/multimetro.jpeg}
  \caption{Modelo dos dois multímetros digitais utilizados para medir
  simultaneamente as tensões nos enrolamentos primário e secundário.}
  \label{fig:multimetro}
\end{figure}

\subsection{Procedimento experimental}

O procedimento foi executado de acordo com o roteiro fornecido durante a aula
\parencite{roteiro}.

Na primeira configuração, a fonte de tensão alternada foi conectada ao
enrolamento de 300 espiras, que passou a constituir o primário. Um multímetro foi
ligado aos terminais desse enrolamento para medir $V_p$, enquanto o segundo foi
conectado ao enrolamento de 1200 espiras para medir simultaneamente $V_s$. A
fonte foi ajustada sucessivamente em dez posições e, em cada uma delas, foram
registradas as duas tensões.

Em seguida, a alimentação foi invertida. A fonte foi conectada ao enrolamento de
1200 espiras, que passou a ser o primário, e o enrolamento de 300 espiras passou
a ser o secundário. Mantendo um multímetro conectado a cada enrolamento, foram
repetidas as medições para as dez posições da fonte.

Para cada linha das duas séries, calculou-se a razão

\begin{equation}
  r_i=\frac{V_{s,i}}{V_{p,i}}.
  \label{eq:razao_experimental}
\end{equation}

A média das dez razões de cada circuito foi determinada por

\begin{equation}
  \bar r=\frac{1}{n}\sum_{i=1}^{n}r_i,
  \label{eq:media}
\end{equation}

e sua dispersão foi expressa pelo desvio padrão amostral

\begin{equation}
  s_r=\sqrt{\frac{\sum_{i=1}^{n}(r_i-\bar r)^2}{n-1}}.
  \label{eq:desvio}
\end{equation}

As dez linhas de cada circuito correspondem a diferentes posições da fonte, e
não a repetições de uma mesma tensão. Por isso, o desvio padrão calculado
caracteriza a dispersão da razão de transformação ao longo das condições
ensaiadas; ele não representa a tolerância dos multímetros. Como não foram
registradas a faixa selecionada em cada instrumento nem suas especificações de
exatidão, não foi possível compor uma incerteza instrumental completa para cada
leitura. A resolução de $0{,}1\,\mathrm{V}$ observada nos registros foi considerada
na discussão qualitativa, mas não foi apresentada como se fosse a tolerância do
equipamento. Essa distinção segue a recomendação de identificar separadamente as
componentes estatísticas e não estatísticas da incerteza
\parencite{jcgm2008}.

A diferença relativa entre a média experimental e o valor teórico foi calculada
por

\begin{equation}
  \delta=\frac{|\bar r-r_{\mathrm{teo}}|}{r_{\mathrm{teo}}}\times100\%.
  \label{eq:diferenca}
\end{equation}

Os cálculos e gráficos foram produzidos em Python com as bibliotecas NumPy e
Matplotlib. O script utilizado é apresentado no Anexo~B.

% =========================================================
% Resultados e discussão
% =========================================================
\section{Resultados e Discussão}

No Circuito 1, o enrolamento de 300 espiras foi utilizado como primário e o de
1200 espiras como secundário. A razão teórica correspondente é

\begin{equation}
  r_{\mathrm{teo},1}=\frac{N_s}{N_p}
  =\frac{1200}{300}=4.
  \label{eq:teorica1}
\end{equation}

Os valores medidos e as razões calculadas são apresentados na
Tabela~\ref{tab:circuito1}.

\begin{table}[H]
  \centering
  \caption{Tensões e razões experimentais do Circuito 1.}
  \label{tab:circuito1}
  \medskip
  \begin{tabular}{cccc}
    \toprule
    Posição da fonte & $V_p$ (V) & $V_s$ (V) & $V_s/V_p$ \\
    \midrule
     1 &  1,9 &  5,2 & 2,7368 \\
     2 &  3,9 & 11,2 & 2,8718 \\
     3 &  6,0 & 17,1 & 2,8500 \\
     4 &  8,2 & 23,1 & 2,8171 \\
     5 & 10,3 & 29,2 & 2,8350 \\
     6 & 12,4 & 35,2 & 2,8387 \\
     7 & 14,5 & 41,2 & 2,8414 \\
     8 & 16,6 & 47,1 & 2,8373 \\
     9 & 18,7 & 52,8 & 2,8235 \\
    10 & 20,8 & 58,6 & 2,8173 \\
    \midrule
    Média & 11,33 & 32,07 & 2,8269 \\
    Desvio padrão de $V_s/V_p$ & -- & -- & 0,0356 \\
    \bottomrule
  \end{tabular}
\end{table}

É importante observar que os valores manuscritos na última coluna da folha de
laboratório correspondem a $V_p/V_s$, apesar de o cabeçalho indicar $V_s/V_p$.
Por exemplo, na primeira linha, $1{,}9/5{,}2=0{,}365$, aproximadamente o valor
anotado, enquanto a razão solicitada é $5{,}2/1{,}9=2{,}7368$. A
Tabela~\ref{tab:circuito1} apresenta todas as razões na ordem indicada pelo
roteiro e pela Equação~\eqref{eq:razao_tensoes}.

Para o Circuito 1, obteve-se a razão adimensional\footnote{Como $V_s$ e $V_p$ são expressos na mesma unidade, o volt se cancela no quociente $V_s/V_p$. Consequentemente, tanto a razão média $\bar r$ quanto seu desvio padrão $s_r$ são adimensionais.}


\begin{equation}
  \boxed{\bar r_1\pm s_{r,1}=2{,}8269\pm0{,}0356}.
\end{equation}

A diferença relativa em relação à razão ideal foi

\begin{equation}
  \delta_1=\frac{|2{,}8269-4|}{4}\times100\%=29{,}33\%.
\end{equation}

Embora \(V_s\) tenha permanecido maior que \(V_p\) em todas as posições, confirmando o funcionamento elevador, a inclinação experimental, aproximadamente igual a \(2{,}83\), ficou abaixo do valor teórico \(4\). Como as razões apresentaram pequena dispersão e permaneceram aproximadamente constantes, a diferença possui caráter predominantemente sistemático. A hipótese mais provável é o acoplamento magnético incompleto entre os enrolamentos, possivelmente agravado por pequenos entreferros nas junções do núcleo ferromagnético desmontável. Além disso, quando a bobina de 300 espiras foi utilizada como primário, sua menor indutância pode ter produzido uma corrente de magnetização mais elevada, aumentando as quedas de tensão na resistência do enrolamento e na impedância interna da fonte. Perdas por histerese e correntes parasitas também podem ter contribuído indiretamente, embora não pareçam suficientes, isoladamente, para explicar a diferença observada. Erros instrumentais são menos prováveis como causa principal, pois seria necessária uma diferença de calibração muito elevada entre os multímetros. Também não se pode excluir uma ligação em uma derivação diferente da nominal. Como não foram medidas a corrente de magnetização, a frequência, as resistências dos enrolamentos nem o fluxo magnético, não é possível quantificar separadamente cada uma dessas contribuições.


No Circuito 2, o enrolamento de 1200 espiras foi alimentado e o de 300 espiras
foi usado como secundário. A razão teórica é

\begin{equation}
  r_{\mathrm{teo},2}=\frac{N_s}{N_p}
  =\frac{300}{1200}=0{,}25.
  \label{eq:teorica2}
\end{equation}

Os resultados dessa configuração são apresentados na
Tabela~\ref{tab:circuito2}.

\begin{table}[H]
  \centering
  \caption{Tensões e razões experimentais do Circuito 2.}
  \label{tab:circuito2}
  \medskip
  \begin{tabular}{cccc}
    \toprule
    Posição da fonte & $V_p$ (V) & $V_s$ (V) & $V_s/V_p$ \\
    \midrule
     1 &  2,1 & 0,4 & 0,1905 \\
     2 &  4,5 & 1,0 & 0,2222 \\
     3 &  7,0 & 1,7 & 0,2429 \\
     4 &  9,4 & 2,3 & 0,2447 \\
     5 & 11,9 & 3,0 & 0,2521 \\
     6 & 14,3 & 3,6 & 0,2517 \\
     7 & 16,8 & 4,3 & 0,2560 \\
     8 & 19,2 & 4,9 & 0,2552 \\
     9 & 21,7 & 5,6 & 0,2581 \\
    10 & 24,2 & 6,3 & 0,2603 \\
    \midrule
    Média & 13,11 & 3,31 & 0,2434 \\
    Desvio padrão de $V_s/V_p$ & -- & -- & 0,0216 \\
    \bottomrule
  \end{tabular}
\end{table}

Para o Circuito 2, obteve-se

\begin{equation}
  \boxed{\bar r_2\pm s_{r,2}=0{,}2434\pm0{,}0216}.
\end{equation}

A diferença relativa em relação ao modelo ideal foi

\begin{equation}
  \delta_2=\frac{|0{,}2434-0{,}25|}{0{,}25}\times100\%=2{,}65\%.
\end{equation}

Nesse circuito, o valor experimental ficou próximo da previsão teórica e
$V_s<V_p$ em todas as posições, caracterizando corretamente o funcionamento
abaixador. As duas primeiras razões foram menores que as demais. Nessas posições,
as tensões secundárias foram apenas $0{,}4\,\mathrm{V}$ e $1{,}0\,\mathrm{V}$;
portanto, a resolução de $0{,}1\,\mathrm{V}$ indicada nas anotações representa
uma fração maior da leitura e exerce maior influência sobre o quociente. A partir
da quinta posição, as razões ficaram entre 0,2517 e 0,2603, muito próximas de
0,25.

A Figura~\ref{fig:tensoes} apresenta $V_s$ em função de $V_p$ para as duas
configurações. Nos dois casos, os pontos formaram tendências aproximadamente
lineares. No Circuito 1, entretanto, observa-se claramente a separação entre os
dados e a reta ideal de inclinação 4; no Circuito 2, as curvas experimental e
ideal permanecem próximas.

\begin{figure}[H]
  \centering
  \imagemrelatorio[width=.98\linewidth]{Figs/tensoes_transformador.png}
  \caption{Tensão secundária em função da tensão primária. As linhas tracejadas
  representam o modelo ideal e as linhas contínuas empregam a razão média
  experimental.}
  \label{fig:tensoes}
\end{figure}

A comparação direta entre as razões teóricas e experimentais é mostrada na
Figura~\ref{fig:comparacao}. As barras de erro representam um desvio padrão das
dez razões medidas, expressando sua dispersão e não uma incerteza instrumental
completa.

\begin{figure}[H]
  \centering
  \imagemrelatorio[width=.78\linewidth]{Figs/comparacao_razoes.png}
  \caption{Comparação das razões teóricas com as médias experimentais. As barras
  de erro correspondem ao desvio padrão amostral.}
  \label{fig:comparacao}
\end{figure}

% =========================================================
% Conclusão
% =========================================================
\section{Conclusão}

O experimento permitiu verificar a relação entre as tensões dos enrolamentos de
um transformador e investigar o efeito da inversão da alimentação. Quando a
bobina de 300 espiras foi utilizada como primário e a de 1200 espiras como
secundário, o sistema funcionou como elevador de tensão. O resultado foi
$V_s/V_p=2{,}8269\pm0{,}0356$, expresso como média $\pm$ desvio padrão, valor
29,33\% inferior à razão
ideal igual a 4.

Após a inversão, com a bobina de 1200 espiras no primário e a de 300 espiras no
secundário, o sistema funcionou como abaixador. Nesse caso, obteve-se
$V_s/V_p=0{,}2434\pm0{,}0216$, também como média $\pm$ desvio padrão,
apresentando diferença relativa
de apenas 2,65\% em relação ao valor teórico de 0,25.

A diferença sistemática observada na configuração elevadora mostrou que o transformador não se comportou como um dispositivo ideal nas condições do ensaio. O resultado pode ser atribuído principalmente ao acoplamento magnético incompleto, possivelmente associado a entreferros no núcleo, à maior corrente de magnetização no enrolamento de 300 espiras e às perdas do conjunto. A pequena dispersão das razões indica que os instrumentos dificilmente constituíram a causa principal da diferença. Dessa forma, o objetivo experimental foi atingido, pois foram confirmadas a proporcionalidade aproximada entre \(V_s\) e \(V_p\) e as funções elevadora e abaixadora do transformador, mantendo-se a necessária distinção entre o modelo ideal e o comportamento real do equipamento.


\newpage

% =========================================================
% Referências
% =========================================================
\printbibliography[heading=bibintoc,title={Referências}]

% =========================================================
% Anexos
% =========================================================
\section*{Anexo A -- Registros do experimento}

\begin{figure}[H]
  \centering
  \imagemrelatorio[width=.93\linewidth]{Figs/dados_laboratorio.jpeg}
  \caption{Folha de dados preenchida durante o experimento.}
  \label{fig:dados_laboratorio}
\end{figure}

\clearpage

\begin{figure}[H]
  \centering
  \imagemrelatorio[width=.72\linewidth]{Figs/quadro_relacoes.jpeg}
  \caption{Relações entre tensões, correntes e números de espiras. Os números de espiras seguem a formulação acima.}
  \label{fig:quadro}
\end{figure}

\clearpage

\section*{Anexo B -- Script Python para o tratamento dos dados}

O Código~\ref{cod:analise} apresenta o script empregado para calcular as razões,
as médias, os desvios padrão amostrais e as diferenças relativas, além de gerar
as Figuras~\ref{fig:tensoes} e~\ref{fig:comparacao}.

\lstinputlisting[
  caption={Script utilizado no tratamento dos dados experimentais.},
  label={cod:analise}
]{analise_transformador.py}

\end{document}
